% =====================================================================
% Mục 5 — Bàn luận  (BẢN TIẾNG VIỆT)
%
% Bản dịch đối ứng 1-1 với sections/05_discussion.tex. Mọi con số phải giống
% bản tiếng Anh. Mục này không đưa ra con số mới nào ngoài Mục 4.
%
% Độ phân tán luôn là độ lệch chuẩn trên ba seed, không bao giờ là ước lượng
% khoảng.
%
% DẤU THẬP PHÂN: giữ dấu chấm như bản tiếng Anh, để phép kiểm tra trùng khớp
% tập con số giữa hai bản còn hiệu lực.
% =====================================================================

\section{Bàn luận}
\label{vi:sec:discussion}

\subsection{Thứ hạng kiến trúc và biến động giữa các seed}
\label{vi:sec:discussion:ranking}

Benchmark tách ra thành hai phát hiện có độ vững chắc rất khác nhau, và không
nên báo cáo chung với cùng một mức tự tin.

Phát hiện vững chắc là ResNet-50 không cạnh tranh được ở bài toán này. Mức
$92.356 \pm 1.022$\,\% của nó thấp hơn ConvNeXt-Tiny khoảng bảy điểm, một
khoảng cách gấp nhiều lần độ lệch chuẩn giữa các seed của chính nó, và nó gây
ra 122 trong 204 lỗi của cả năm backbone trên mọi seed
(Mục~\ref{vi:sec:results:perspecies}). Khoảng cách này còn mang tính cấu trúc
chứ không khuếch tán: nó tập trung ở một số loài cụ thể mà bốn kiến trúc kia
gần như hoàn hảo.

Phát hiện mong manh là thứ tự giữa bốn mô hình còn lại. ConvNeXt-Tiny hơn
Swin-Tiny 0.438 điểm độ chính xác trong khi độ lệch chuẩn của hai bên là 0.354
và 0.581; MobileNetV3-Large kém Swin-Tiny 0.752 điểm với độ lệch chuẩn 0.319.
Các khoảng cách này cùng bậc độ lớn với biến động giữa các seed, nên với ba
seed thì thứ tự không thể coi là đã xác lập, và một bộ ba seed khác hoàn toàn
có thể đảo lại. Khẳng định bảo vệ được là ConvNeXt-Tiny, Swin-Tiny và
MobileNetV3-Large đều đạt tới một điểm vận hành tương đương, và việc chỉ ra
một kiến trúc tốt nhất trong số đó sẽ đòi hỏi nhiều seed hơn nghiên cứu này đã
chạy.

Kết quả của MobileNetV3-Large đáng được nhấn mạnh vì một lý do khác. Nó đạt
trong phạm vi 1.19 điểm độ chính xác so với mô hình tốt nhất, với 4.22\,M tham
số so với 27.83\,M, và 0.22\,G MACs so với 4.46\,G
(Bảng~\ref{tab:efficiency}), tức khoảng một phần sáu số tham số và một phần
hai mươi khối lượng tính toán. Với kịch bản triển khai tại hiện trường (động
lực chính của việc định danh gỗ bằng ảnh mặt cắt), đánh đổi đó quan trọng hơn
thứ hạng ở đỉnh bảng. Đây chỉ là phép so sánh giữa độ chính xác và độ phức tạp
tĩnh. Độ trễ suy luận không được báo cáo ở đây, với lý do nêu ở
Mục~\ref{vi:sec:results:efficiency}.

\subsection{Cấu trúc nhầm lẫn nội họ: Cupressaceae so với Pinaceae}
\label{vi:sec:discussion:cupressaceae}

Độ khó ở cấp họ không tỷ lệ với số loài trong họ. Cupressaceae là họ lớn nhất
của bộ dữ liệu với tám loài, nhưng lại có tỷ lệ lỗi thấp nhất trong năm họ, chỉ
0.78\,\%.

Điều đáng phân tích nằm ở chỗ khác, là vị trí của lỗi so với ranh giới họ. Các
số đếm dưới đây cộng theo từng phần tử trên 15 lần chạy chính của registry,
cùng cách gộp dùng cho Hình~\ref{fig:confusion}. Chỉ 12 trong 30 lỗi của
Cupressaceae là nhầm lẫn nội họ, và gần như toàn bộ số đó
thuộc một cặp duy nhất, \textit{Cunninghamia lanceolata} và
\textit{Callitris columellaris}, với 6 lần theo chiều thứ nhất và 5 lần theo
chiều ngược lại. Cặp này gần như cân bằng. Pinaceae cho một bức ảnh trái ngược:
75 trong 78 lỗi của họ này là nội họ, và riêng 61 lỗi dồn vào đúng một chiều,
\textit{Pinus krempfii} bị dự đoán thành \textit{Pinus dalatensis}.

Hai dạng cấu trúc đó gợi ý hai cơ chế khác nhau. Một cặp nhỏ và đối xứng phù
hợp với cách đọc rằng hai loài thực sự khó phân biệt ở mức phân giải hiện có,
khi mô hình nhầm theo cả hai chiều với tần suất tương đương. Một cặp lớn và bất
đối xứng thì không đọc được như vậy, vì nếu chỉ là khó phân biệt thuần túy thì
lỗi phải xuất hiện theo cả hai chiều. Sự bất đối xứng 61 so với một gợi ý mô
hình có thiên lệch hệ thống về phía một lớp, và thiên lệch đó không giảm khi
năng lực mô hình tăng.

Hệ quả cho việc đọc kết quả cấp họ là tỷ lệ lỗi của một họ không đủ để mô tả độ
khó của nó. Cupressaceae và Pinaceae có tỷ lệ lỗi lần lượt là 0.78\,\% và
2.83\,\%, chênh nhau chưa tới bốn lần, nhưng cấu trúc lỗi bên dưới hai con số
đó khác nhau về bản chất.

\subsection{Hiệu năng trên các loài đặc hữu Việt Nam}
\label{vi:sec:discussion:endemic}

Bốn loài trong bộ dữ liệu là loài đặc hữu của Việt Nam
[CẦN NGUỒN: tài liệu phân loại học xác nhận tính đặc hữu của bốn loài này].
Đây là nhóm loài mà một công cụ định danh triển khai trong nước sẽ gặp thường
xuyên nhất, nên hiệu năng riêng trên nhóm này đáng được tách ra xem xét.

Kết quả không đồng nhất trong nhóm. Mọi giá trị F1 trích dẫn ở đây đều lấy từ
Bảng~\ref{tab:perspecies_full}. Dưới ConvNeXt-Tiny,
\textit{Xanthocyparis vietnamensis} và \textit{Pinus cernua} đều đạt
$100.000 \pm 0.000$\,\% F1 trên cả ba seed. Hai loài còn lại thấp hơn rõ rệt,
với \textit{Pinus dalatensis} đạt $96.559 \pm 3.295$\,\% và
\textit{Pinus krempfii} đạt $96.573 \pm 3.495$\,\%. Chúng cũng là hai loài mang
độ lệch chuẩn giữa các seed lớn nhất trong cột này.

Như vậy khó khăn còn lại của mô hình tốt nhất tập trung vào đúng hai loài đặc
hữu cùng thuộc chi \textit{Pinus}, chứ không gắn với tính đặc hữu nói chung.
Cách đọc này được củng cố bởi hai loài phân bố tại Việt Nam nhưng không thuộc
chi \textit{Pinus}: \textit{Calocedrus rupestris} và
\textit{Fokienia hodginsii} đạt $100.000 \pm 0.000$\,\% F1 trên cả năm
backbone, tức không một kiến trúc nào mắc lỗi trên hai loài này. Yếu tố quyết
định vì vậy là quan hệ phân loại học gần giữa hai loài trong cùng một chi,
không phải nguồn gốc địa lý của mẫu vật.

Cần đặt kết quả này cạnh giới hạn nêu ở Mục~\ref{vi:sec:discussion:shortcut}.
Điểm F1 cao trên các loài đặc hữu được đo trên cùng bộ sưu tập và cùng điều
kiện chụp với tập huấn luyện, nên chúng chưa nói được gì về hiệu năng trên mẫu
vật thu thập độc lập.

\subsection{Mô hình đang dùng gì: bằng chứng phản bác cách đọc theo giải phẫu}
\label{vi:sec:discussion:shortcut}

Giới hạn hệ trọng nhất của benchmark này không nằm ở độ lớn các khoảng cách độ
chính xác, mà ở \emph{ý nghĩa} của chính các con số độ chính xác. Các phép thăm
dò ở Mục~\ref{vi:sec:results:convergence} cho thấy một Swin-Tiny huấn luyện
trên đầu vào $16\times16$ đạt 94.361\,\% độ chính xác, và một mô hình huấn
luyện trên ảnh làm mờ nặng đạt 95.677\,\%, so với 99.060\,\% của mức tham
chiếu độ phân giải đầy đủ. Cả hai phép biến đổi đều phá hủy cấu trúc giải phẫu
mà việc định danh gỗ bằng mặt cắt dựa vào, và cả hai đều giữ lại trên 95\,\%
hiệu năng của ảnh đầy đủ.

Các con số thăm dò dẫn ở đây nằm ở Bảng~\ref{tab:probe}, còn mức tham chiếu
toàn ảnh nằm ở Bảng~\ref{tab:task_b}. Cách giải thích đầu tiên, cho rằng mô
hình chỉ đọc màu sắc, không được ủng hộ: phép thăm dò chỉ dùng histogram đạt
37.970\,\%. Màu sắc có mang tín hiệu
thật, cao hơn hẳn mức ngẫu nhiên 5.882\,\%, nhưng nó chỉ chiếm chưa đến một nửa
những gì hai phép thăm dò có thông tin không gian đạt được. Vậy bằng chứng
quyết định mang tính không gian nhưng ở tần số thấp: một thứ sống sót qua việc
lấy mẫu giảm mười sáu lần. Độ dốc dài theo tổng thể, vân thớ thô, phân bố độ
sáng toàn cục và các đặc trưng thu ảnh gắn với từng mẫu vật đều phù hợp với
quan sát này, và các thí nghiệm hiện tại không tách được chúng ra.

Còn lại hai cách đọc, với hệ quả rất khác nhau. Cách đọc lành tính là các loài
cây lá kim thực sự khác nhau về tính chất quang học tổng thể ở thang đo này,
khi đó mô hình đã học được một dấu hiệu có thật nhưng thô hơn dự định, và
benchmark đo một bài toán chính đáng dù hẹp hơn. Cách đọc bất lợi là nội dung
tần số thấp mang \emph{danh tính mẫu vật} (như ánh sáng, thuốc nhuộm, cách gắn
mẫu, thiết lập máy ảnh) và rò rỉ mẫu vật còn sót lại
(Mục~\ref{vi:sec:discussion:originkey}) cho phép khai thác điều đó. Dữ liệu
hiện có không tách được hai cách đọc. Để phân biệt chúng cần đánh giá trên ảnh
được thu thập độc lập, thứ mà nghiên cứu này không có
(Mục~\ref{vi:sec:discussion:external}).

Hệ quả cho việc báo cáo là trực tiếp. Các mức độ chính xác ở tiêu đề nên đọc
như độ chính xác trên bài toán này và bộ dữ liệu này, không phải như bằng chứng
rằng mô hình định danh gỗ bằng giải phẫu. Chính quy trình thăm dò của nhóm
cũng ghi nhận benchmark chưa được phát hành trên cơ sở đó. Điều này không làm
mất hiệu lực phép so sánh kiến trúc ở
Mục~\ref{vi:sec:discussion:ranking}: cả năm backbone đều được huấn luyện và
đánh giá trong điều kiện như nhau, nên thứ tự tương đối giữa chúng vẫn là một
phép so sánh hợp lệ. Nó chỉ giới hạn điều mà các con số \emph{tuyệt đối} được
phép khẳng định.

Hai quan sát nữa rút ra từ các phép thăm dò. Thứ tự tỷ lệ giữ lại (làm mờ cao
hơn lấy mẫu giảm) gợi ý rằng nội dung hữu ích nằm ở thang đo thô hơn mức mà
phép lấy mẫu $16\times16$ giữ được, và đây là một dự đoán kiểm chứng được bằng
một phép quét độ phân giải (thí nghiệm đã được cấu hình cho nghiên cứu này
nhưng chưa bao giờ được chạy, xem
Mục~\ref{vi:sec:discussion:scope}). Và các phép thăm dò chỉ huấn luyện 20
epoch so với 50 của các lần chạy chính, nên con số tuyệt đối của chúng là cận
dưới của những gì đầu vào đã bị phá hủy có thể hỗ trợ, điều này làm mạnh
thêm chứ không làm yếu lập luận.

\subsection{Xoay ảnh không cải thiện đồng đều}
\label{vi:sec:discussion:rotation}

Xoay rời rạc $90^\circ$ xuất phát từ nhận xét rằng một mặt cắt ngang không có
hướng chuẩn tắc, nên các phép xoay bội của $90^\circ$ lẽ ra phải bảo toàn mọi
đặc trưng chẩn đoán. Dữ liệu không ủng hộ lợi ích được kỳ vọng đó.

Hiệu ứng đổi dấu giữa các kiến trúc: dương với ResNet-50 và EfficientNet-B3,
đúng bằng không với MobileNetV3-Large, âm với cả ConvNeXt-Tiny và Swin-Tiny
(Bảng~\ref{tab:ablation}). Một phép tăng cường thực sự mã hóa được một tính
chất của hình học thu ảnh thì lẽ ra phải giúp ích, hoặc ít nhất không gây hại,
bất kể backbone nào. Việc hai mô hình mạnh nhất lại chính là hai mô hình giảm
là ngược với quy luật mà một lợi ích bất biến thật sự sẽ tạo ra.

Còn lại hai cách đọc mà nghiên cứu này không tách được. Phép tăng cường có thể
là thừa, bởi cấu hình chính đã áp dụng xoay liên tục $\pm30^\circ$ cùng lật
ngang và lật dọc, nên phần lớn tính bất biến theo hướng có thể đã được phủ,
khiến xoay rời rạc $90^\circ$ đánh đổi đa dạng màu mà không được gì. Hoặc các
chênh lệch đơn giản là tạp âm seed: bốn trong năm độ lớn nằm trong khoảng 1.6
lần độ lệch chuẩn giữa các seed tương ứng, và với một seed cho mỗi ablation thì
không có cách nào phân biệt một hiệu ứng thật nhỏ với một lần bốc may mắn. Dù
theo hướng nào, kết luận thực tế vẫn thế: không nên lấy xoay rời rạc
$90^\circ$ làm mặc định cho bài toán này dựa trên những lần chạy ở đây.

Label smoothing hành xử khác và đáng tách riêng. Bỏ nó làm giảm độ chính xác ở
cả năm backbone, không ngoại lệ nào. Không chênh lệch đơn lẻ nào lớn, và ba
trong năm nằm dưới độ lệch chuẩn giữa các seed tương ứng, nên độ lớn hiệu ứng
chưa được xác lập. Dẫu vậy, sự nhất quán về \emph{hướng} trên năm kiến trúc
độc lập là quy luật duy nhất trong bảng ablation không giống tạp âm. Nâng label
smoothing lên 0.2 không cho sự nhất quán như vậy. Đây là bằng chứng yếu rằng
giá trị mặc định 0.1 là một lựa chọn hợp lý, và \emph{không} phải bằng chứng
rằng 0.1 là tối ưu.

\subsection{Thiết kế ablation một seed là giới hạn ràng buộc}
\label{vi:sec:discussion:oneseed}

Mọi ablation trong nghiên cứu này chỉ huấn luyện với seed 42. Đây là giới hạn
phương pháp hệ trọng nhất của các kết quả ablation, và nó không sửa được bằng
phân tích.

Các lần chạy chính cho thấy độ lệch chuẩn giữa các seed ở bài toán này nằm
trong khoảng 0.319 đến 1.022 điểm độ chính xác tùy backbone
(Bảng~\ref{tab:task_b}). Trong 15 chênh lệch ablation ở
Bảng~\ref{tab:ablation}, có 9 nhỏ hơn độ lệch chuẩn của chính backbone tương
ứng, và chỉ một vượt qua nó hơn khoảng 2.4 lần. Vậy một chênh lệch một seed có
độ lớn vài phần của một điểm gần như không mang thông tin gì về hiệu ứng bên
dưới: nhiễu đo và hiệu ứng đang đo có cùng một cỡ.

Điều này có một hệ quả cụ thể về cách đọc bảng ablation. Không nên diễn giải
từng ô riêng lẻ. Thứ đọc được là quy luật trên cả hàng năm backbone, thể hiện qua một
dấu nhất quán trên năm kiến trúc khởi tạo độc lập, như thấy ở label smoothing
0.0, khó giải thích bằng tạp âm seed hơn bất kỳ ô đơn lẻ nào, ngay cả khi từng
ô một mình không mang thông tin. Đó là lý do
Mục~\ref{vi:sec:results:ablation} báo cáo sự nhất quán về hướng chứ không phải
độ lớn hiệu ứng, và không ablation nào được trình bày như một cải tiến.

Lặp lại mỗi ablation trên đúng ba seed đã dùng cho các lần chạy chính sẽ tốn
thêm 30 lần huấn luyện và sẽ làm cho các phép so sánh này diễn giải được. Đây
là mở rộng có giá trị cao nhất cho nghiên cứu này.

\subsection{Ảnh hưởng của cách xây dựng khóa nguồn}
\label{vi:sec:discussion:originkey}

Phân hoạch Task~B gom ảnh theo mẫu vật sao cho mọi lát cắt từ cùng một mẫu gỗ
vật lý đều nằm trong một tập. Điều này quan trọng vì bộ dữ liệu gốc bị cắt rất
nhiều: 5{,}830 ảnh sinh ra từ 557 mẫu vật, và 88.3\,\% số ảnh đến từ những mẫu
vật đóng góp từ 24 lát cắt trở lên. Với phép chia ở cấp từng ảnh, một mẫu vật
gồm 30 lát cắt có xác suất khoảng $6.9\times10^{-16}$ để không rơi vào tập
huấn luyện, nên gần như mọi mẫu vật lớn đều xuất hiện trong tập train, và ảnh
test trở thành những lát cắt khác của mẫu đã quen thay vì mẫu mới. Phân hoạch
Task~B cấp-ảnh trước đây đúng là như vậy: 97.62\,\% số ảnh test của nó đến từ
những mẫu vật đã có mặt trong tập train. Độ chính xác đo dưới phân hoạch đó mô
tả khả năng nhận lại lát cắt, không phải khả năng khái quát hóa theo loài.

Cách gom theo mẫu vật dùng ở đây loại bỏ đúng chế độ thất bại đó, nhưng bảo đảm
mà nó cung cấp bị giới hạn bởi \emph{cách xây dựng khóa mẫu vật}. Khóa được suy
ra từ tên tệp: bỏ chỉ số lát cắt ở cuối và thêm tên lớp ở đầu, nên
\texttt{3351.\ Araucaria araucana.1.(12).png} trở thành
\texttt{Araucaria araucana::3351.\ Araucaria araucana.1.}. Đó là một quy ước
đặt tên, không phải định danh mẫu vật được ghi lúc thu thập. Từ đó có hai hệ
quả khác nhau, và chúng không đối xứng.

Trường hợp gom thừa, khi hai mẫu vật khác nhau về vật lý nhưng trùng gốc tên tệp nên bị
nhập làm một khóa, chỉ làm mất tính linh hoạt khi chia. Đây là hướng bảo thủ
và không làm thổi phồng kết quả.

Gom thiếu mới là rủi ro thật. Nếu các lát cắt của cùng một mẫu vật được lưu
dưới những gốc tên khác nhau, ví dụ qua nhiều phiên quét riêng biệt hoặc sau
một bước đổi tên, thì khóa sẽ gán cho chúng những nguồn khác nhau và phép gom
\emph{âm thầm} thất bại trong việc tách chúng ra. Không có gì trong tên tệp
phát hiện được điều này, và không có định danh mẫu vật độc lập nào trong bộ dữ
liệu để đối chiếu.

Một cuộc kiểm toán riêng về khóa này cho thấy độ lớn của mối lo ngại đó. Gom
lại 552 khóa của tập 17 lớp theo một quy tắc ứng viên thô hơn cho ra 325 nhóm,
tức thu gọn 41.1\,\%: theo cách đọc đó, bốn trong mười khóa là mảnh của một đơn
vị vật lý lớn hơn chứ không phải là đơn vị độc lập. Kiểm toán cũng tìm thấy 86
mã số mẫu bị dùng lại giữa các loài, nên phần số trong tên tệp tự nó không phải
là một định danh mẫu vật. Kết luận chính thức của kiểm toán là \textbf{chưa thể
xác định} khóa này định danh mẫu vật vật lý hay chỉ là một ảnh/góc chụp, và
rằng cấu trúc tên tệp không được thay thế cho bước xác minh đối chiếu với sổ
mẫu của phòng thí nghiệm; bước xác minh đó vẫn chưa hoàn tất. Do vậy, độ mạnh
của khẳng định tách mẫu vật phụ thuộc vào kỷ luật đặt tên của bộ sưu tập gốc,
điều mà nghiên cứu này không tự kiểm chứng được.

Hai lưu ý nữa ảnh hưởng đến những gì kiểm tra được từ các artifact đã công bố.
Thứ nhất, khoảng 19\,\% ảnh test mang tên tệp đã đánh số lại trong phân hoạch
triển khai, nên không thể tái lập khóa nguồn của chúng về sau; trong 81\,\% tái
lập được, không mẫu vật nào xuất hiện ở cả validation lẫn test. Thứ hai,
registry chỉ công bố dự đoán cho tập validation và test, nên việc tách
train--test (yếu tố thực sự quan trọng) không thể kiểm chứng lại từ các
artifact đã phát hành và phụ thuộc vào chính cách xây dựng phân hoạch. Công bố
kèm bảng ánh xạ mẫu-vật sang tập sẽ đóng được lỗ hổng này.

\subsection{Chưa có kiểm chứng ngoài hoặc tại hiện trường}
\label{vi:sec:discussion:external}

Toàn bộ kết quả đến từ một bộ dữ liệu duy nhất, được chia nội bộ. Các mô hình
chưa bao giờ được đánh giá trên gỗ được chụp bởi một người khác, bằng thiết bị
khác, dưới ánh sáng khác, hoặc tại một địa điểm khác. Độ chính xác trên
99\,\% nên đọc là độ chính xác trên những mẫu vật giữ lại \emph{của bộ sưu tập
này}, không phải là ước lượng hiệu năng tại hiện trường.

Giới hạn này chặt hơn về bề ngoài, vì một lý do đặc thù của nghiên cứu này.
Nhầm lẫn duy nhất còn sót lại ở mô hình mạnh nhất là giữa hai loài
\textit{Pinus}, và các lỗi tách ResNet-50 khỏi phần còn lại tập trung ở một vài
loài. Cả hai quy luật đều là thuộc tính của bộ dữ liệu này ở quy mô hiện tại
(từ 24 đến 39 ảnh test mỗi loài theo Bảng~\ref{tab:perspecies_full}, rút ra từ
một số lượng mẫu vật hạn chế mỗi
loài). Một bộ sưu tập khác hoàn toàn có thể cho ra một tập loài khó khác mà
không mâu thuẫn với bất kỳ điều gì báo cáo ở đây.

Không có gì trong nghiên cứu này ủng hộ một khẳng định về mức độ sẵn sàng triển
khai. Bước tiếp theo phù hợp là đánh giá trên một tập mẫu vật được thu thập độc
lập, lý tưởng là chụp trong những điều kiện mô hình chưa từng thấy, trước khi
cân nhắc bất kỳ sử dụng vận hành nào.

\subsection{Quyết định về phạm vi: Task A và ablation độ phân giải}
\label{vi:sec:discussion:scope}

Hai thí nghiệm có trong kế hoạch dự án nhưng không được báo cáo ở
Mục~\ref{vi:sec:results}, với lý do khác nhau.

\paragraph{Task A (phân loại cấp họ).} Nhiệm vụ cấp họ dùng phép chia phân tầng
theo loài, thao tác ở cấp từng ảnh, đúng cơ chế đã tạo ra 97.62\,\% độ phủ
mẫu vật ở phân hoạch Task~B trước đây
(Mục~\ref{vi:sec:discussion:originkey}), áp lên cùng bộ dữ liệu bị cắt nhiều
đó. Kết quả của nó do vậy chịu cùng một rò rỉ, và các điểm gần như tuyệt đối ở
cấp họ, kể cả một kiến trúc đạt 100.000\,\% trên cả ba seed, không thể đọc như
bằng chứng về khả năng khái quát hóa sang mẫu vật chưa thấy. Báo cáo Task~A
cạnh một Task~B đã gom theo mẫu vật sẽ đặt hai giao thức đánh giá không so sánh
được cạnh nhau trong cùng một mục kết quả, và mời gọi đúng cách đọc sai mà phép
gom sinh ra để ngăn chặn. Vì vậy Task~A bị \emph{loại} chứ không báo cáo kèm
cảnh báo. Khôi phục nó sẽ đòi hỏi dựng lại phép chia cấp họ theo mẫu vật và
huấn luyện lại cả 15 lần chạy; ở cấp họ điều này ít rủi ro thiếu lớp hơn ở cấp
loài, vì ba loài bị loại khỏi Task~B thuộc những họ vẫn còn nhiều mẫu vật khác.
Cũng lưu ý rằng nhiệm vụ cấp họ được định nghĩa trên bốn họ, gộp
\textit{Araucaria araucana} vào Pinaceae, trong khi cách nhóm dùng ở
Mục~\ref{vi:sec:results} giữ Araucariaceae riêng; hai sơ đồ này phải được thống
nhất trước khi có thể báo cáo hai nhiệm vụ cạnh nhau.

\paragraph{Độ phân giải đầu vào.} Một ablation độ phân giải $384$ pixel đã được
cấu hình nhưng chưa bao giờ được huấn luyện; không có lần chạy nào cho nó trong
registry. Độ phân giải đầu vào do vậy chưa được kiểm chứng trong nghiên cứu
này. Đây là một khoảng trống thật sự chứ không phải một sự loại bỏ: độ phân
giải có thể tương tác với đúng cái khó khăn còn sót lại ở mô hình tốt nhất, vì
việc phân biệt \textit{P.\ krempfii} với \textit{P.\ dalatensis} có thể phụ
thuộc vào chi tiết giải phẫu tinh mà đầu vào $224$ pixel không phân giải được.
Kiểm chứng điều này là phần bổ tự nhiên cho ablation nhiều seed đề xuất ở
Mục~\ref{vi:sec:discussion:oneseed}.
